% Part: history
% Chapter: biographies
% Section: bertrand-russell

\documentclass[../../../include/open-logic-section]{subfiles}

\begin{document}

\olfileid[hi]{his}{bio}{rus} 

\olsection{बर्ट्रैंड रसेल}

\olphoto{russell-bertrand}{बर्ट्रैंड रसेल}

बर्ट्रैंड रसेल को आधुनिक विश्लेषणात्मक दर्शन के संस्थापकों में से एक माना
जाता है। 18 मई 1872 को जन्मे रसेल दर्शन और तर्कशास्त्र में अपने कार्य के
लिए ही नहीं जाने जाते थे; उन्होंने विविध विषयों पर अनेक लोकप्रिय पुस्तकें भी
लिखीं। वे जीवन भर प्रबल राजनीतिक कार्यकर्ता भी रहे।

रसेल का जन्म वेल्स के मॉनमथशायर में ट्रेलेख में हुआ। उनके माता-पिता ब्रिटिश
अभिजात वर्ग के सदस्य थे। वे स्वतंत्र विचारक थे और उस समय बोस्टन के उग्र
सुधारवादियों से भी उनकी मित्रता थी। दुर्भाग्य से रसेल के बचपन में ही उनके
माता-पिता की मृत्यु हो गई और उन्हें दादा-दादी के साथ रहने भेज दिया गया। वहाँ
उनका धार्मिक पालन-पोषण हुआ (जिससे उनके माता-पिता हर कीमत पर बचना चाहते थे)।
नैतिकता के सभी मामलों में उनकी दादी अत्यंत कठोर थीं। किशोरावस्था में निजी
शिक्षकों ने अधिकतर उन्हें घर पर पढ़ाया।

विश्लेषणात्मक दर्शन और विशेषतः तर्कशास्त्र पर रसेल का प्रभाव अत्यंत व्यापक
है। उन्होंने ट्रिनिटी कॉलेज, कैम्ब्रिज में गणित और दर्शन पढ़ा, जहाँ गणितज्ञ
और दार्शनिक अल्फ्रेड नॉर्थ व्हाइटहेड ने उन्हें प्रभावित किया। 1910 में रसेल
और व्हाइटहेड ने \emph{Principia Mathematica} का पहला खंड प्रकाशित किया,
जिसमें उन्होंने इस मत का समर्थन किया कि गणित को तर्कशास्त्र में घटाया जा
सकता है। आगे उन्होंने सैकड़ों पुस्तकें, निबंध और राजनीतिक पुस्तिकाएँ
प्रकाशित कीं। 1950 में उन्हें साहित्य का नोबेल पुरस्कार मिला।

रसेल राजनीति और सामाजिक सक्रियता में गहराई से संलग्न थे। प्रथम विश्व
युद्ध के दौरान शांतिवादी गतिविधियों और विरोध के कारण उन्हें गिरफ़्तार करके
छह महीने के लिए कारागार भेजा गया। कारागार में वे लिख-पढ़ सके और उनका कहना
था कि यह अनुभव ``काफ़ी सुखद'' था। वे जीवन भर शांतिवादी रहे और 1961 में
परमाणु निरस्त्रीकरण की एक सभा में भाग लेने पर फिर कारावास हुआ। वे 1948 की
एक विमान दुर्घटना में भी बचे, जिसमें केवल धूम्रपान खंड में बैठे लोग जीवित
बचे थे। इसलिए रसेल कहते थे कि उनका जीवन धूम्रपान के कारण बचा। रसेल ने चार
विवाह किए, किंतु विवाहेतर संबंधों के लिए भी उनकी ख्याति थी। उनका निधन 2
फ़रवरी 1970 को वेल्स के पेनरिनडाइद्राइथ में 97 वर्ष की आयु में हुआ।

\begin{reading}
रसेल ने 1872--1967 के अपने जीवन को समेटती तीन भागों की आत्मकथा लिखी
\citep{Russell1967,Russell1968,Russell1969}। मैकमास्टर विश्वविद्यालय का
बर्ट्रैंड रसेल अनुसंधान केंद्र, बर्ट्रैंड रसेल अभिलेखागार का घर है। उनकी
संकलित रचनाओं के खंडों (खोज-योग्य अनुक्रमणिकाओं सहित) और अभिलेखीय
परियोजनाओं की जानकारी के लिए उनकी वेबसाइट \citet{Duncan2015} देखें। रसेल
का शोधपत्र \emph{On Denoting} \citep{Russell1905}, बीसवीं शताब्दी के
विश्लेषणात्मक दर्शन की एक उत्कृष्ट कृति है।

रसेल पर \emph{Stanford Encyclopedia of Philosophy} की प्रविष्टि
\citep{Irvine2015} में इच्छा और राजनीतिक सिद्धांत पर रसेल के वक्तव्य की
ध्वनि-कतरनें हैं। रसेल के कई वीडियो साक्षात्कार ऑनलाइन उपलब्ध हैं। उदाहरणतः
उन्हें धूम्रपान और विमान दुर्घटना में होने की चर्चा करते देखने के लिए
\citet{RussellND} देखें। रसेल की कुछ रचनाएँ, जिनमें उनकी
\emph{Introduction to Mathematical Philosophy} भी है, \citet{LibriVoxND}
पर निःशुल्क श्रव्य-पुस्तकों के रूप में उपलब्ध हैं।
\end{reading}

\end{document}
